\section[gRPC-интерфейс для высоконагруженных расчётов]{Проектирование gRPC-интерфейса для высоконагруженных расчётов}

Помимо пользовательских приложений, в проекте реализован сервисный контур, предназначенный для интеграции вычислительного ядра во внешнюю инфраструктуру.  
В отличие от настольных сценариев, здесь основной режим работы связан с пакетной обработкой большого числа точек и включением расчётов в состав других программных систем.  

Сервисный слой реализован отдельным компонентом \path{if97_calculator_service}.  
Транспортным протоколом выбран gRPC, который обеспечивает строгий контракт на уровне протокола и эффективную сериализацию данных~\cite{grpc_docs}.  

Ключевым методом API является двунаправленный потоковый вызов \texttt{If97Service/CalculatePt}.  
Метод ориентирован на обработку батчей однотипных запросов:
независимые вычисления $pT$ выполняются параллельно,
а результаты возвращаются по мере готовности.  

Формат сообщений построен по схеме «массивы полей» (Struct-of-Arrays, SoA).  
Смысл SoA состоит в том, что данные передаются не как список объектов «строка таблицы», а как набор примитивных массивов одинаковой длины.  
Например, запрос содержит идентификатор батча и массивы входных параметров:
\begin{itemize}
\item сообщение \texttt{PtBatchRequest} содержит идентификатор батча и массивы входных параметров $p[]$ и $t[]$;  
\item сообщение \texttt{PtBatchResponse} содержит идентификатор батча, массив $h[]$, массив статусов $status[]$ и поля диагностики уровня батча \texttt{batch\_error\_*}.  
\end{itemize}

В такой схеме \texttt{status[i]} соответствует строке входного батча
и позволяет локализовать ошибку без остановки всего потока.
Поля \texttt{batch\_error\_*} используются для ошибок уровня батча
(например, несовпадение длин входных массивов).

На рисунке~\ref{fig:grpc-sequence-if97-service} показан полный маршрут
пакетного запроса в сервисе:
от приёма HTTP/2-потока на стороне \texttt{tonic}
до передачи данных в \texttt{if97\_core}
и формирования \texttt{PtBatchResponse}.
Схема подчёркивает, что сетевой контур,
асинхронная обработка и вычислительное ядро разделены,
но при этом работают как единый поток обработки батча.

\begin{figure}[htbp]
    \centering
    % \textwidth заставит диаграмму растянуться или сжаться ровно под ширину текста
    % ! сохранит пропорции (высота подстроится автоматически)
    \resizebox{\linewidth}{!}{
        \begin{tikzpicture}[
            >=Latex,
            participant/.style={
                draw=black,
                fill=gray!15,
                minimum width=3.2cm,
                minimum height=0.9cm,
                align=center,
                font=\small
            },
            lifeline/.style={
                draw=black,
                densely dashed,
                line width=0.5pt
            },
            message/.style={
                draw=black,
                line width=0.9pt,
                -Latex
            },
            act/.style={
                draw=black,
                fill=gray!25,
                line width=0.5pt
            }
        ]

        % Participants
        \node[participant] (client) at (0,0) {Клиент};
        \node[participant] (http2)  at (4.7,0) {HTTP/2\\(Сетевой стек)};
        \node[participant] (tonic)  at (9.4,0) {Tonic\\(gRPC сервер)};
        \node[participant] (core)   at (14.1,0) {if97\_core\\(Вычислительное ядро)};

        % Lifelines
        \draw[lifeline] (0,-0.45) -- (0,-13.7);
        \draw[lifeline] (4.7,-0.45) -- (4.7,-13.7);
        \draw[lifeline] (9.4,-0.45) -- (9.4,-13.7);
        \draw[lifeline] (14.1,-0.45) -- (14.1,-13.7);

        % Activations
        \draw[act] (9.22,-3.35) rectangle (9.58,-11.55);
        \draw[act] (13.92,-6.55) rectangle (14.28,-9.75);

        % Messages
        \draw[message] (0,-1.6) -- node[above, font=\scriptsize, align=center] {1. Отправка пакетного запроса\\Protobuf-сообщение \texttt{PtBatchRequest}} (4.7,-1.6);

        \draw[message] (4.7,-3.0) -- node[above, font=\scriptsize, align=center] {2. Мультиплексированный поток\\\texttt{CalculatePt} по HTTP/2 (порт 50051)} (9.4,-3.0);

        \node[font=\scriptsize, align=left, anchor=west] at (10.95,-4.95) {3. Десериализация\\и валидация батча};
        \draw[message] (9.4,-4.5) -- (10.65,-4.5) -- (10.65,-5.3) -- (9.4,-5.3);

        \draw[message] (9.4,-6.4) -- node[above, font=\scriptsize, align=center] {4. Асинхронная передача данных\\в \texttt{Dispatcher} / Rayon ThreadPool} (14.1,-6.4);

        \node[font=\scriptsize, align=left, anchor=west] at (15.55,-7.65) {5. Параллельный расчет\\(маршрутизация и математика)};
        \draw[message] (14.1,-7.2) -- (15.35,-7.2) -- (15.35,-8.0) -- (14.1,-8.0);

        \draw[message] (14.1,-9.2) -- node[above, font=\scriptsize, align=center] {6. Возврат массива результатов\\\texttt{h[]} и \texttt{status[]}} (9.4,-9.2);

        \draw[message] (9.4,-10.8) -- node[above, font=\scriptsize, align=center] {7. Формирование Protobuf-ответа\\\texttt{PtBatchResponse}} (4.7,-10.8);

        \draw[message] (4.7,-12.2) -- node[above, font=\scriptsize, align=center] {8. Возврат ответа} (0,-12.2);

        \end{tikzpicture}
    }
    
    \caption{Диаграмма последовательности обработки пакетного gRPC-запроса в сервисе \texttt{if97\_calculator\_service}: от приема HTTP/2 streaming-запроса до параллельного вычисления в \texttt{if97\_core} и формирования Protobuf-ответа.}
    \label{fig:grpc-sequence-if97-service}
\end{figure}

Преимущества данного подхода:
\begin{itemize}
\item уменьшаются накладные расходы на сериализацию и разбор protobuf-сообщений;  
\item упрощается эффективная параллельная обработка данных внутри вычислительного пула;  
\item снижается количество промежуточных аллокаций при подготовке данных;  
\item сохраняется возможность построчно возвращать статус (корректно рассчитано/ошибка) и диагностическую информацию.  
\end{itemize}

\section[Параллелизм и асинхронная обработка запросов]{Модель параллелизма и асинхронная обработка запросов средствами Tokio}

Высокопроизводительная обработка запросов в сервисе требует явного разделения сетевого ввода-вывода и вычислений.  
Термодинамические расчёты по IF97 являются вычислительно интенсивными задачами, поэтому их выполнение в сетевой среде исполнения без изоляции приводит к деградации задержек и пропускной способности.  

В сервисе применена двухконтурная модель исполнения:
\begin{itemize}
\item модуль \texttt{tokio} отвечает за gRPC-среду исполнения, сетевые соединения и асинхронный ввод-вывод;  
\item модуль \texttt{rayon} выполняет вычисления в пуле фиксированного размера.  
\end{itemize}
Разделение асинхронного ввода-вывода и вычислительного пула
соответствует используемому технологическому стеку
~\cite{tokio_spawn_blocking_docs,rayon_docs}.  

Поток обработки батча включает следующие этапы.  
Сначала сервис принимает очередной фрагмент входного потока, выполняет проверку формата и ограничений.  
Далее батч передаётся в вычислительный пул, где вычисления выполняются параллельно по строкам.  
После завершения вычислений результаты возвращаются в поток ответов
через очередь \texttt{mpsc} с учётом ограничения давления очередей.

Для устойчивости под нагрузкой реализован механизм ограничения давления очередей на уровне количества батчей, одновременно находящихся в обработке.  
Ограничения задаются переменными окружения и используются как предохранители от неконтролируемого роста очередей и потребления памяти:
\begin{itemize}
\item переменная \path{IF97_MAX_IN_FLIGHT_PER_CONN} задаёт лимит на одну клиентскую сессию;  
\item переменная \path{IF97_MAX_IN_FLIGHT_GLOBAL} задаёт глобальный лимит на весь процесс сервиса;  
\item переменная \path{IF97_MAX_BATCH_LEN} задаёт ограничение длины батча;  
\item переменная \path{IF97_GRPC_MAX_MSG_BYTES} задаёт ограничение размера gRPC-сообщения.  
\end{itemize}

Для эксплуатации также важны проверка работоспособности и корректное завершение.  
Сервис поднимает \texttt{tonic-health} и использует статусы \texttt{Serving/NotServing}
~\cite{tonic_health_docs}.  
При остановке включается режим штатного опустошения очереди (drain):
процесс переводится в \texttt{NotServing},
перестаёт принимать новый трафик
и завершает активные вычисления в пределах \texttt{IF97\_DRAIN\_SECS}.  
